\subsection{Purpose}

Account and access-management procedures ensure that access to systems, applications, networks, devices, and data is limited to authorized users and services based on business need.

Access \must{} be granted, used, reviewed, modified, and removed in a controlled and documented manner.

\subsection{Account Requirements}

Each person \must{} use an individually assigned account for normal system and application access.

Individual accounts \must{} be used to support accountability, auditing, access review, and incident investigation.

Accounts \must{} be associated with a responsible person, role, service, or system.

Accounts \must{} have an owner and a defined purpose. Unowned, unidentified, or unnecessary accounts \must{} be disabled or removed.

Users \mustnot{} share passwords, authentication devices, recovery codes, API keys, or other authentication information.

Shared accounts \should{} not be used. When a shared account is required for technical or operational reasons:
\begin{itemize}
    \item The account \must{} have a documented business purpose
    \item Use of the account \must{} be approved by the system owner
    \item Access \must{} be limited to authorized personnel
    \item Activity \must{} be logged where technically practicable
    \item The account credentials \must{} be changed when an authorized user leaves the role
    \item The account \must{} be reviewed periodically
\end{itemize}

\subsection{Access Requests and Approval}

Access \must{} be approved before it is provided.

Access requests \must{} identify:
\begin{itemize}
    \item The person or service requesting access
    \item The system, application, network, or data requested
    \item The requested role or permission level
    \item The business purpose
    \item The requested start date and, where applicable, end date
    \item The approving manager or system owner
\end{itemize}

Access \must{} be based on the principles of least privilege and need to know.

Users \must{} receive only the access required to perform their assigned duties.

Access to sensitive, critical, administrative, or financial systems \must{} require approval from the appropriate system or data owner.

High-risk access \should{} require approval from more than one authorized person where separation of duties is reasonably practicable.

\subsection{Privileged Access}

Privileged access includes administrative, root, domain, database, cloud-management, security-management, and other access capable of changing system configurations, permissions, security controls, or sensitive data.

Privileged access \must{} be restricted to authorized administrators and other personnel with a documented business requirement.

Administrators \must{} use separate privileged accounts for administrative activities and ordinary accounts for routine work.

Privileged accounts \mustnot{} be used for routine email, web browsing, or other ordinary user activities unless technically necessary and approved.

Privileged access \should{} use multifactor authentication and approved secure management methods.

Privileged access \must{} be logged where technically practicable. Administrative logs \must{} record enough information to identify the account, action, affected system, and date and time.

Privileged access \must{} be reviewed periodically and after changes in responsibilities.

\subsection{Authentication}

Systems must use approved authentication controls appropriate to the sensitivity and risk of the system.

Multifactor authentication \must{} be enabled for:
\begin{itemize}
    \item Administrative and privileged accounts
    \item Remote access
    \item Cloud-management interfaces
    \item Access to sensitive or regulated information
    \item Systems where multifactor authentication is supported and required by the system owner
\end{itemize}

Authentication credentials \must{} be protected from unauthorized disclosure.

Passwords \must{} comply with the password-management requirements in this document.

Systems \should{} prevent the reuse of recently used passwords, limit repeated failed authentication attempts, and require reauthentication after an appropriate period of inactivity.

Recovery codes, backup authentication methods, and emergency credentials \must{} be protected to the same extent as passwords.

\subsection{Account Provisioning}

New accounts \must{} be created only after the required approval has been obtained.

Before an account is created, the administrator \should{} verify:
\begin{itemize}
    \item The user's identity or service ownership
    \item The business purpose
    \item The requested access level
    \item The approving authority
    \item Any start date, expiration date, or access limitations
\end{itemize}

Accounts for temporary workers, contractors, vendors, interns, and other non-permanent users \must{} have an expiration date or scheduled review date.

Service accounts \must{} have a documented owner, purpose, required permissions, and credential-management method.

Service accounts \mustnot{} be used for interactive logins unless technically required and approved.

\subsection{Changes to Access}

Changes to a user's role, responsibilities, employment status, contract, or system assignment \must{} result in a review of that user's access.

Access that is no longer required \must{} be removed promptly.

Administrators \must{} not increase permissions merely to resolve an operational issue without documenting and approving the change.

Temporary elevated access \must{} have a defined duration and \must{} be removed when the approved period ends.

\subsection{Account Termination}

User access \must{} be disabled or removed when employment, contract, volunteer service, or other authorization ends.

Termination procedures \must{} address, as applicable:
\begin{itemize}
    \item Network and application accounts
    \item Privileged and administrative accounts
    \item Remote-access accounts
    \item Email and collaboration accounts
    \item Cloud and vendor accounts
    \item Physical access credentials
    \item Company-owned devices
    \item Tokens, keys, certificates, and authentication devices
    \item Shared-account passwords known to the departing person
\end{itemize}

For planned terminations, access removal \must{} occur no later than the end of the person's authorization.

For urgent or involuntary terminations, access \must{} be removed or suspended immediately when directed by authorized management.

\subsection{Access Reviews}

Access rights \must{} be reviewed periodically by the appropriate manager, system owner, or data owner.

Access reviews \should{} confirm:
\begin{itemize}
    \item The user or service still requires access
    \item The assigned permissions are appropriate
    \item Privileged access remains necessary
    \item Former users and inactive accounts have been removed
    \item Temporary access has expired
    \item Service accounts have current owners
    \item Shared accounts remain necessary and controlled
\end{itemize}

Critical systems, privileged accounts, and access to sensitive information \should{} be reviewed at least quarterly.

Other accounts and access rights \should{} be reviewed at least annually or more frequently based on risk.

Review results \must{} be documented, including the reviewer, date, systems reviewed, identified issues, and corrective actions.

\subsection{Inactive and Dormant Accounts}

Accounts that are no longer required \must{} be disabled or removed.

Accounts that have not been used for a defined period \should{} be disabled automatically where technically practicable.

Dormant accounts \must{} be reviewed before being reactivated.

Default, test, demonstration, and guest accounts \must{} be disabled or removed unless there is a documented business requirement.

\subsection{Emergency and Break-Glass Accounts}

Emergency or break-glass accounts may be used only when normal administrative access is unavailable or insufficient to respond to an urgent operational or security need.

Emergency accounts \must{}:
\begin{itemize}
    \item Have a documented purpose and owner
    \item Be protected using strong credentials
    \item Be restricted to authorized personnel
    \item Be monitored and logged where technically practicable
    \item Be tested periodically
    \item Be reviewed after every use
    \item Have their credentials changed after use when appropriate
\end{itemize}

Use of an emergency account \must{} be documented, including the reason, user, time, actions performed, and outcome.

\subsection{Third-Party Access}

Third-party access \must{} be approved before it is enabled.

Third-party access \must{} be limited to:
\begin{itemize}
    \item The systems and data required for the approved work
    \item The minimum permissions required
    \item The approved time period
    \item Approved connection methods
\end{itemize}

Third-party access \must{} use individual accounts where technically practicable.

Vendor and contractor accounts \must{} be disabled or removed when the work is complete or the relationship ends.

Third-party access \must{} be reviewed periodically and after significant changes to the contract, service, system, or risk.

\subsection{Account and Access Records}

Account and access records \must{} be maintained for systems that require formal access management.

Records \should{} include:
\begin{itemize}
    \item Account name and type
    \item Account owner
    \item User, service, or vendor associated with the account
    \item System or application
    \item Assigned role and permissions
    \item Approval information
    \item Creation date
    \item Expiration or review date
    \item Last access review
    \item Account status
\end{itemize}

Access records \must{} be protected from unauthorized modification and disclosure.